\documentclass[a4paper,10pt]{article}

\usepackage[top=0.38in, bottom=0.38in, left=0.55in, right=0.55in]{geometry}
\usepackage{enumitem}
\usepackage{hyperref}
\usepackage{titlesec}
\usepackage{array}
\usepackage[T1]{fontenc}
\usepackage{lmodern}
\usepackage{microtype}
\hypersetup{colorlinks=true, urlcolor=black, linkcolor=black}
\pagestyle{empty}
\titleformat{\section}{\bfseries\normalsize}{}{0em}{}[\vspace{-4pt}\rule{\linewidth}{0.5pt}]
\titlespacing{\section}{0pt}{4pt}{1pt}
\setlist[itemize]{leftmargin=1.1em,itemsep=0pt,parsep=0pt,topsep=1pt,label=\textbullet}

\begin{document}

\begin{center}
    {\Large\bfseries WENTWORTH LIU}\\[2pt]
    {\small Beijing, China \;$\cdot$\; \href{mailto:gzliu0705@gmail.com}{gzliu0705@gmail.com} \;$\cdot$\; \href{https://www.linkedin.com/in/wentworth-liu-4a4a02413/}{linkedin.com/in/wentworth-liu-4a4a02413}}\\[2pt]
    {\small Data Scientist \;$\cdot$\; Data Automation \& Model Workflows \;$\cdot$\; Scientific Computing \;$\cdot$\; Operational Analytics}
\end{center}

\vspace{-3pt}

\section{SUMMARY}
Data Scientist with 2+ years of experience building automated, data-driven workflows for complex engineering and operational systems. Strong in Python and SQL, with experience in multi-source data integration, validation, transformation, configurable feature systems, and reusable analytical workflows for million-scale EV fleets at Li Auto. Brings a Computational Materials Science background and a careful, automate-first approach to improving data quality, consistency, and decision-support outputs.

\section{EXPERIENCE}

\noindent\textbf{Li Auto} \hfill Beijing, China\\
Data Scientist, Battery Intelligence Department \hfill Jul 2024 -- Present
\begin{itemize}
    \item \textbf{Automated Data Transformation \& Model Workflows.} Built a configurable feature system for multi-dimensional fleet telemetry using custom event definitions, multi-operator aggregations, and cross-platform normalisation; deployed reusable workflows to an internal ML platform and improved feature-iteration efficiency by \textasciitilde85\%.

    \item \textbf{Data Validation \& Analytical Outputs.} Built an explainable risk-ranking pipeline across multiple battery platforms, producing vehicle-level risk scores and feature-level drivers; achieved 93\% Recall@2\% on retrospectively validated high-risk cases while limiting review to the highest-risk 2\% of vehicles.

    \item \textbf{Scenario Analysis \& Decision Support.} Used quasi-causal analysis on million-scale telemetry to identify high-risk SOH-degradation cohorts and actionable drivers; productised intervention analysis across multiple vehicle models to support consistent scenario assessment and targeted decisions.

    \item \textbf{Workflow Automation.} Built a hybrid rule-based and statistical classifier for battery-quality work orders across 30+ fault categories, combining text matching, TF-IDF, and Random Forest; achieved Weighted F1 > 0.90 and used confidence-based routing to focus review on ambiguous cases.
\end{itemize}

\vspace{2pt}
\noindent\textbf{ByteDance} \hfill Beijing, China\\
Data Analytics Intern, Douyin Market Strategy \hfill May -- Sep 2023
\begin{itemize}
    \item \textbf{Multi-source Integration \& Data Monitoring.} Integrated consumption, dissemination, and search signals across multiple data platforms into an IP popularity measurement system; added automated pipeline monitoring to identify upstream data instability and support reliable recurring analysis.

    \item \textbf{Metric Design \& Analytical Reporting.} Redesigned performance metrics and built a multi-dimensional content taxonomy for Qishui Music; used funnel segmentation and audience-saturation modelling to diagnose conversion decline and inform content and resource-allocation decisions.
\end{itemize}

\vspace{2pt}
\noindent\textbf{Dongwu Securities} \hfill China\\
Equity Research Intern, Mechanical Sector \hfill Dec 2021 -- Mar 2022
\begin{itemize}
    \item \textbf{Quantitative Analysis \& Reporting.} Combined company fundamentals, competitive analysis, and market dynamics with stock-price modelling to deliver structured investment recommendations for Harmonic Drive Systems.
\end{itemize}

\section{EDUCATION}

\noindent\textbf{Nanjing University} \hfill Sep 2021 -- Jun 2024\\
MSc Computational Materials Science
\begin{itemize}
    \item Developed Skeaf, a 3D Fermi-surface visualisation tool combining HDBSCAN clustering with quantum-oscillation data; first author, Physical Review B (2024); co-author, Physical Review B (2024) and Physical Review Materials (2023); Patent: CN115795266A.
\end{itemize}

\vspace{2pt}
\noindent\textbf{Nanjing University} \hfill Sep 2017 -- Jul 2021\\
BSc Materials Science \& Engineering (Outstanding Graduate) \;$\cdot$\; GPA: 4.6/5.0 \;$\cdot$\; Top 1\% of cohort
\begin{itemize}
    \item National Scholarship (top 0.15\%) \;$\cdot$\; MCM/ICM 2020 Honorable Mention (Team Captain) \;$\cdot$\; Sakura Science Exchange Programme, Tohoku University, Japan (5 selected university-wide).
\end{itemize}

\section{SKILLS}

\begin{tabular}{@{}p{3.0cm}p{12.5cm}@{}}
\textbf{Programming} & Python \textperiodcentered{} SQL/HiveSQL \textperiodcentered{} Pandas \textperiodcentered{} NumPy \textperiodcentered{} C++ \textperiodcentered{} MATLAB \textperiodcentered{} Git \\[1pt]
\textbf{Data Automation} & ETL Pipelines \textperiodcentered{} Data Validation \textperiodcentered{} Data Transformation \textperiodcentered{} Multi-source Data Integration \textperiodcentered{} Data Standardisation \textperiodcentered{} Automated Monitoring \textperiodcentered{} Workflow Automation \textperiodcentered{} Reporting \textperiodcentered{} Visualisation \\[1pt]
\textbf{Analytics \& Delivery} & Statistical Modelling \textperiodcentered{} Predictive Modelling \textperiodcentered{} Machine Learning \textperiodcentered{} Feature Engineering \textperiodcentered{} Time Series \textperiodcentered{} Causal Inference \textperiodcentered{} Data Productisation \textperiodcentered{} GCP \textperiodcentered{} Docker \textperiodcentered{} CI/CD \textperiodcentered{} Spark \textperiodcentered{} Hadoop \textperiodcentered{} Linux \\[1pt]
\textbf{Domain \& Lang.} & Scientific Computing \textperiodcentered{} Engineering Data \textperiodcentered{} EV Analytics \textperiodcentered{} English (TOEFL 104) \textperiodcentered{} Mandarin/Cantonese (Native) \textperiodcentered{} Japanese (JLPT N1) \\
\end{tabular}

\end{document}
